Dependent type theory can be viewed as both a programming language and a foundation of mathematics.
In particular, the Curry-Howard correspondence allows us to use dependently typed programming languages such as \texttt{Agda}, \texttt{Dedukti}, \texttt{Lean}, and \texttt{Rocq} as proof assistants.
To ensure consistency, these languages should only accept total functions, so they put bounds on recursive calls.
For example, if the function has an inductive domain (freely generated by constructors), then we know that the function terminates if it is only called recursively on a smaller input.
Dually, if the function has a coinductive codomain (cofreely generated by destructors), then we know that the function is productive if its output is larger than that of its recursive calls.
In both of these cases, we know that the function is total; however, totality can also be the result of a more intricate mix of induction and coinduction.

\subparagraph{Syntactic Totality.}
Totality is usually enforced via syntactic checks for termination and productivity.
By default, \texttt{Rocq} and \texttt{Lean} check for structural (co)recursion \cite{goguen_eliminating_2006,cockx_eliminating_2016,thibodeau_intensional_2020}, while \texttt{Agda} and \texttt{Dedukti} use the more permissive size-change termination principle \cite{size-change}.
Both conditions rely on a syntactic comparison of the structural size of the input/output of the original function with the structural size of the input/output of the recursive call. 
However, programs with more complex patterns of (co)recursion are not easily defined in a way that satisfies these syntactic conditions.
Moreover, syntactic methods are sensitive to slight reformulations of programs, and do not work well with higher-order programs \cite{Abel-inflationary}.

\subparagraph{Type-based Totality.}
This work is concerned with a type-based approach to enforcing totality: sized types.
Sized types were originally introduced by Hughes, Pareto, and Sabry \cite{hughes} as an alternative approach to both termination and productivity checking, in which these properties are guaranteed by the typing derivation instead of an inspection of syntax.
In this approach, both inductive and coinductive types can be annotated with additional size information.

We can think of sizes as ordinals, and they come equipped with a well-founded ordering.
In particular, for every size \(i\) and inductive type \(\sf{Ind}\) we have an ordinal approximation \(\sf{Ind}^i\), which can be seen informally as the subtype of \(\sf{Ind}\) consisting of terms that can be formed within \(i\) constructor steps.
Dually, for every size \(i\) and coinductive type \(\sf{CoInd}\) we have an ordinal approximation \(\sf{CoInd}^i\), which can be seen informally as the quotient of \(\sf{CoInd}\), where terms are identified if we cannot distinguish them within \(i\) destructor steps.

By defining functions on these annotated versions of the (co)inductive types, we introduce information about the sizes of the input/output directly in the type of the function.
This allows for greater compositionality of functions, since size-information is retained across function calls.
Moreover, it allows us to reduce both termination and productivity of recursive functions to well-founded induction on the size parameter.

Sized types, and similar approaches, have been proposed in different forms to aid  termination/productivity checking \cite{xi_sizes_2002, blanqui_sizes_2004, barthe_sizes_2004, Abel2006TypebasedTA, abel_2008}. 
The combination of sized types with dependent type theory has been studied as well \cite{Barthe, sacchini_2013, sacchini_2014, NbE, Chan_2022}.
Sized types have been implemented in \texttt{Agda} \cite{PAR-10:MiniAgda_Integrating_Sized_Dependent}, and there have been informative practical uses, including: encoding regular languages coalgebraically \cite{abel_2016}, up-to techniques for bisimilarity \cite{up-to-sized-types}, and constructing quotient-inductive types \cite{constructing-qits}.

\subparagraph*{Consistency.}
Importantly, \texttt{Agda} does not allow pattern matching on sizes: they should have no influence on computation, and are only used to aid the termination checker.
However, there are still significant limitations to both the current implementation of sized types, and to other proposals.
One design question is how the full (co)inductive type \((\sf{Co})\sf{Ind}\) should relate to its ordinal approximations \((\sf{Co})\sf{Ind}^i\).

\texttt{Agda} introduces a special largest size $\infty$, which can be thought of as a closure ordinal: an ordinal large enough that the approximation has reached the full (co)inductive type, so that we can define \((\sf{Co})\sf{Ind}\coloneqq (\sf{Co})\sf{Ind}^\infty\).
However, this approach leads to inconsistencies: \texttt{Agda} proves \(\infty<\infty\), while any strict well-order \(<\) can be proven to be irreflexive \cite{agda-issue-3026}.

We employ a different way of recovering the full (co)inductive type from its ordinal approximations: using parametric quantifiers \cite{reynolds1983types, parametric-quantifier, reldtt}.
The intuition is that the quantifiers \(\exists/\forall\) are versions of \(\Sigma/\Pi\) that restrict information: the type \(\exists i.\,A(i)\) contains dependent pairs where the first component is ignored when checking equality, while the universal quantifier \(\forall i.\,A(i)\) contains dependent functions that act uniformly on any input.
With these quantifiers, we can define the full inductive type as \(\sf{Ind}\coloneqq\exists i.\,\sf{Ind}^i\), and the full coinductive type as \(\sf{CoInd}\coloneqq\forall i.\,\sf{CoInd}^i\).
This follows to our intuition: a term of an inductive type should be bounded by some possibly infinite height (where the bound is kept abstract), whereas a term of a coinductive type should be observable up to any depth (where the depth does not influence observations). 
Using \(\Sigma/\Pi\) instead of \(\exists/\forall\) would make \((\sf{Co})\sf{Ind}\) too large, so our parametric quantifiers have the same effect as forbidding pattern matching on sizes.

To our knowledge, using parametric quantifiers for sizes was first described by Vezzosi \cite{licentiate}, although it was framed in the context of guarded type theory \cite{nakano,prod-coprogramming} to extend its applicability from coinductive to inductive types.
It was later adapted by Nuyts, Vezzosi and Devriese in work on dependent type theory with parametric quantifiers \cite{parametric-quantifier, reldtt}, where this parametric approach to sized types was described as a possible application. In their work, they also give a way to construct (co)inductive types: they construct all coinductive types but only the finitely-branching inductive types; this is because the type of sizes is interpreted as the set of natural numbers in their model.
The interpretation of sizes as a larger ordinal has been explored \cite{Abel2006TypebasedTA, Chan_2022}, but not in combination with parametricity for sizes, which is central to our construction of (co)inductive types.

\subparagraph*{Contribution.}
We show how all (co)inductive types can be encoded using sizes and parametric quantifiers.
Concretely, we extend a version of Martin-L\"of type theory --- containing an empty type, unit type, boolean type, \(\Sigma\)-types, \(\Pi\)-types, \(=\)-types, one universe, function extensionality, and no primitive (co)inductive types --- with two additional features: a type of sizes, and parametric quantifiers \(\exists\) and \(\forall\) over sizes.
In our formulation, parametricity is internalised via the addition of axioms, and quantifying over small types will produce a small type although the type of sizes is large.

We show how well-founded recursion on sizes can be used to construct ordinal approximations for any (co)inductive type.
Then, we use the parametric quantifiers to construct the full (co)inductive types: we construct both the initial algebra and the final coalgebra for any polynomial endofunctor.
Our definitions do not use the uniqueness of identity proofs (UIP) axiom, and are therefore compatible with homotopy type theory \cite{hottbook}.

We show the consistency of our theory by constructing a realisability model based on Hyland's effective topos \cite{hyland_effective_1982}: small types are interpreted as partial equivalence relations, and large types as assemblies.
This lets us interpret the type of sizes as an uncountable ordinal, which is needed to validate the axioms for the parametric quantifiers.
This models a stronger type theory: one that is also impredicative and extensional and therefore satisfies UIP.

Importantly, we do not use impredicativity and extensionality in the syntax when constructing the ordinal approximations and full (co)inductive types.
So, our constructions can be carried out in a weak theory, while our model validates a strong theory.

\subparagraph*{Related work.}
Beyond the work that we already mentioned it is worth noting the similarity between this approach to sized types and guarded type theory \cite{nakano}, as this was also the context in which this parametric approach to sizes was described \cite{licentiate}.
Guarded type theory introduces a modality $\RHD$, called the later modality.
This modality allows for the construction of guarded recursive types, in which the recursive variable only occurs under the guarding modality.
However, the guarding modality by itself is not sufficient to construct coinductive types.
Instead, the type system can be extended with so-called clock variables \cite{prod-coprogramming}, which represent different time streams and can be used to define coinductive types via universal quantification over clocks.
This construction requires quantification over clocks to be parametric, which is often expressed in the form of additional axioms \cite{mogelberg,guardedtypes16}.
This is very similar to our approach, as we encode coinductive types using parametric universal quantification over sizes. 
However, we also encode inductive types using existential quantification, whereas guarded type theory focuses on constructing coinductive types.
The relation between sized types and guarded type theory has been studied, and it was shown that guarded recursion can be simulated in a simple type theory using sized types \cite{guarded-via-sized}.

Another method for guaranteeing termination is via user-defined well-founded orderings \cite{Leroy2024}, which is possible in \texttt{Lean}, or using even weaker orderings \cite{qwf-order}.
\texttt{Rocq} also has the \texttt{equations} package \cite{hutchison_equations_2010,mangin2015equations,sozeau2019equations} that extends the termination checker at the cost of losing computational behaviour.
Although more flexible, these methods require significant manual work from the user.

\subparagraph*{Structure of the paper.}
\autoref{sec:type-theory} defines the theory in which we work, introducing the type of sizes and parametric quantification. In \autoref{sec:encoding}, we develop the encoding the approximations and for the full (co)inductive types by quantification over sizes. The consistency of the theory is justified in \autoref{sec:model}, which describes the main features of the realisability model. Finally, a summary and perspective on future work are given in \autoref{sec:conclusion}.